\documentclass[a4j,12pt,fleqn,dvipdfmx]{jsarticle}
\usepackage{lib/Quiz}
\usepackage{Setup}
\input{lib/DRAFT_SWITCH}

\definecolor{Red}{RGB}{255,0,0}
\definecolor{Blue}{RGB}{0,0,255}
\definecolor{Green}{RGB}{0,166,93}
\def\Red{\color{Red}}
\def\Blue{\color{Blue}}
\def\Green{\color{Green}}

\gdef\reportinstruction{□《提出方法》 まずは，しっかりとした答案（図示を含む）を，ノートに作成してください．得られた選択肢を，C-Learning 科目名 → 小テスト →
「{\reporttitle}…」より入力・送信してください（ノートの内容は提出しません）．□}

\begin{document}

\begin{Quiz}{1}{小テストタイトル} %{講義回}{小テストタイトル}

\vskip1zh

シリアルナンバー.tex の「シリアルナンバー」は，選択肢をシャッフルする乱数のシードとして使われる．
以下の空欄に相応しい内容を選択肢から選べ.

\paragraph{\Red 重大な注意点}
同じ内容の正答が存在すると，答案が一意でなくなり，不当な減点が起こります．



\paragraph{問1}

次の空欄に相応しい図を選べ．

%%%
%%% \Qbox{配点}{正答}
%%%
\begin{center}
  \Qbox{2}{\includegraphics[scale=0.3]{lib/sample1.eps}}
\end{center}

\paragraph{問2}

次の空欄に相応しい式を選べ．
%%%
%%% amsmath の align, gather 等を使うと連番が狂う．
%%%
\begin{equation}
  I = \Qbox{2}{$\displaystyle \int_{-1}^{1}\sin x\,dx$}
\end{equation}
インライン\Qbox{2}{$\sin x$}でも大丈夫．もちろん\Qbox{2}{文字}でも大丈夫．

%%%
%%% \Qdummy{誤答}
%%%
\Qdummy{\includegraphics[scale=0.3]{lib/sample2.eps}}
\Qdummy{$\displaystyle \int_{-1}^{1}\cos x\,dx$}
\Qdummy{$\tan x$}

\vskip2zh

\end{Quiz}
\end{document}
